\documentclass[UTF8]{ctexart}
\usepackage{xeCJK}
\usepackage{geometry}
\geometry{left=2cm,right=2cm,top=2.5cm,bottom=2.5cm}
\setlength{\headheight}{12.7pt}

\usepackage{titlesec}
\titleformat{\section}
  {\normalfont\large\bfseries}
  {\thesection}
  {1em}
  {}
\titlespacing*{\section}
  {0pt}
  {3.5ex plus 1ex minus .2ex}
  {2.3ex plus .2ex}

\usepackage{amsmath}
\usepackage{amssymb}
\usepackage{amsthm}
\usepackage{enumitem}
\setlist[enumerate]{itemsep=0pt,topsep=0pt,parsep=0pt,leftmargin=0pt,itemindent=2.5em}
\setlist[itemize]{itemsep=0pt,topsep=0pt,parsep=0pt,leftmargin=0pt,itemindent=2.5em}
\usepackage{fancyhdr}
\pagestyle{fancy}
\fancyhf{}
\usepackage{graphicx}
\usepackage{hyperref}
\usepackage{indentfirst}
\usepackage{lmodern}


\begin{document}
\setlength{\abovedisplayskip}{1pt}
\setlength{\belowdisplayskip}{1pt}

\section{函子}

\hypertarget{sec:定义1.1}{}
\noindent\textbf{定义1.1 (协变函子)}

给定\href{01 范畴#nameddest=sec:定义1.1}{范畴}
$\mathbf{C},\mathbf{D}$. 如果$F=(F_O,F_M)$满足
\begin{enumerate}
    \item $F_O:\mathrm{Ob}(\mathbf{C})\to\mathrm{Ob}(\mathbf{D})$,
    $F_M:\mathrm{Mor}(\mathbf{C})\to\mathrm{Mor}(\mathbf{D})$,
    \item 对$\mathbf{C}$中任意的$f:X\to Y$, 在$\mathbf{D}$中有
    $F_M(f):F_O(X)\to F_O(Y)$,
    \item 对$\mathbf{C}$中任意的对象$X$, $F_M(I_X)=I_{F_O(X)}$,
    \item 对$\mathbf{C}$中任意的态射$f,g$, 在有定义的情形下
    $F_M(g\circ f)=F_M(g)\circ F_M(f)$ (即下图是交换图),
    \begin{center}\begin{tikzpicture}[
        box/.style={draw, thick, minimum width=5.5cm, minimum height=2.4cm},
        arrow/.style={-Stealth, thick}]
        \node[box] (Cbox) at (-4,0) {};
        \node[anchor=north west, xshift=5, yshift=-5] at (Cbox.north west)
        {$\mathbf{C}$};
        \node (X) at (-6, -0.6) {$X$};
        \node (Y) at (-4, 0.6) {$Y$};
        \node (Z) at (-2, -0.6) {$Z$};
        \draw[arrow] (X) -- node[left, pos=0.7, xshift=-5] {$f$} (Y);
        \draw[arrow] (Y) -- node[right, pos=0.3, xshift=7] {$g$} (Z);
        \draw[arrow] (X) -- node[below, pos=0.5] {$g\circ f$} (Z);
        \node[box] (Dbox) at (4,0) {};
        \node[anchor=north west, xshift=5, yshift=-5] at (Dbox.north west)
        {$\mathbf{D}$};
        \node (FX) at (2, -0.6) {$F(X)$};
        \node (FY) at (4, 0.6) {$F(Y)$};
        \node (FZ) at (6, -0.6) {$F(Z)$};
        \draw[arrow] (FX) -- node[left, pos=0.7, xshift=-5] {$F(f)$} (FY);
        \draw[arrow] (FY) -- node[right, pos=0.3, xshift=5] {$F(g)$} (FZ);
        \draw[arrow] (FX) -- node[below, pos=0.5] {$F(g\circ f)$} (FZ);
        \draw[arrow, very thick] ([xshift=10]Cbox.east) -- node[above, pos=0.5]
        {$F$} ([xshift=-10]Dbox.west);
    \end{tikzpicture}\end{center}
\end{enumerate}
\vspace{-3pt}
就称$F$是$\mathbf{C}$到$\mathbf{D}$的\textbf{函子}或\textbf{协变函子},
在不产生歧义的情况下记为$F:\mathbf{C}\to\mathbf{D}$.

\noindent{\kaishu 备注.}
在不产生歧义的情况下把$F_O$和$F_M$都统一写成$F$.

\vspace{10pt}

记$\mathbf{C}$到$\mathbf{D}$的全体函子组成的集合为$[\mathbf{C},\mathbf{D}]$.

\vspace{10pt}

\hypertarget{sec:定义1.2}{}
\noindent\textbf{定义1.2 (反变函子)}

给定范畴$\mathbf{C},\mathbf{D}$. 如果$F=(F_O,F_M)$满足
\begin{enumerate}
    \item $F_O:\mathrm{Ob}(\mathbf{C})\to\mathrm{Ob}(\mathbf{D})$,
    $F_M:\mathrm{Mor}(\mathbf{C})\to\mathrm{Mor}(\mathbf{D})$,
    \item 对$\mathbf{C}$中任意的$f:X\to Y$, 在$\mathbf{D}$中有
    $F_M(f):F_O(Y)\to F_O(X)$,
    \item 对$\mathbf{C}$中任意的对象$X$, $F_M(I_X)=I_{F_O(X)}$,
    \item 对$\mathbf{C}$中任意的态射$f,g$, 在有定义的情形下
    $F_M(g\circ f)=F_M(f)\circ F_M(g)$ (即下图是交换图),
    \begin{center}\begin{tikzpicture}[
        box/.style={draw, thick, minimum width=5.5cm, minimum height=2.4cm},
        arrow/.style={-Stealth, thick}]
        \node[box] (Cbox) at (-4,0) {};
        \node[anchor=north west, xshift=5, yshift=-5] at (Cbox.north west)
        {$\mathbf{C}$};
        \node (X) at (-6, -0.6) {$X$};
        \node (Y) at (-4, 0.6) {$Y$};
        \node (Z) at (-2, -0.6) {$Z$};
        \draw[arrow] (X) -- node[left, pos=0.7, xshift=-5] {$f$} (Y);
        \draw[arrow] (Y) -- node[right, pos=0.3, xshift=7] {$g$} (Z);
        \draw[arrow] (X) -- node[below, pos=0.5] {$g\circ f$} (Z);
        \node[box] (Dbox) at (4,0) {};
        \node[anchor=north west, xshift=5, yshift=-5] at (Dbox.north west)
        {$\mathbf{D}$};
        \node (FX) at (2, -0.6) {$F(X)$};
        \node (FY) at (4, 0.6) {$F(Y)$};
        \node (FZ) at (6, -0.6) {$F(Z)$};
        \draw[arrow] (FY) -- node[left, pos=0.3, xshift=-5] {$F(f)$} (FX);
        \draw[arrow] (FZ) -- node[right, pos=0.7, xshift=5] {$F(g)$} (FY);
        \draw[arrow] (FZ) -- node[below, pos=0.5] {$F(g\circ f)$} (FX);
        \draw[arrow, very thick] ([xshift=10]Cbox.east) -- node[above, pos=0.5]
        {$F$} ([xshift=-10]Dbox.west);
    \end{tikzpicture}\end{center}
\end{enumerate}
\vspace{-3pt}
就称$F$是$\mathbf{C}$到$\mathbf{D}$的\textbf{反变函子}.

\vspace{10pt}

$\mathbf{C}\to\mathbf{D}$的反变函子等同于$\mathbf{C}^{\mathrm{op}}\to\mathbf{D}$
(或$\mathbf{C}\to\mathbf{D}^{\mathrm{op}}$)的协变函子.

\vspace{10pt}

\hypertarget{sec:定义1.3}{}
\noindent\textbf{定义1.3 (忠实、全、本质满)}

给定范畴$\mathbf{C},\mathbf{D}$和函子$F:\mathbf{C}\to\mathbf{D}$.
\begin{enumerate}
    \item 如果对$\mathbf{C}$中任意的对象$X,Y$,
    $F_M\!\!\restriction_{\mathbf{C}(X,Y)}:\mathbf{C}(X,Y)\to\mathbf{D}(F(X),F(Y))$
    是单射, 就称$F$是\textbf{忠实}的.
    \item 如果对$\mathbf{C}$中任意的对象$X,Y$,
    $F_M\!\!\restriction_{\mathbf{C}(X,Y)}:\mathbf{C}(X,Y)\to\mathbf{D}(F(X),F(Y))$
    是满射, 就称$F$是\textbf{全}的.
    \item 如果对$\mathbf{D}$中任意的对象$X'$, 存在$\mathbf{C}$中的对象$X$
    使得$F(X)\cong X'$, 就称$F$是\textbf{本质满}的.
\end{enumerate}





\newpage

下面给出一些函子的例子.

\vspace{10pt}

\hypertarget{sec:例1.1}{}
\noindent\textbf{例1.1 (包含函子、恒等函子)}

给定范畴$\mathbf{C}$和它的子范畴$\mathbf{C}'$,
则$\iota:(X\mapsto X,f\mapsto f)$是$\mathbf{C}'$到$\mathbf{C}$的函子,
称之为\textbf{包含函子}.

当$\mathbf{C}'=\mathbf{C}$时, 该$\iota$称为$\mathbf{C}$上的\textbf{恒等函子},
记作$I_{\mathbf{C}}$.

\vspace{10pt}

\hypertarget{sec:定理1.1}{}
\noindent\textbf{定理1.1}

任给范畴$\mathbf{C}$和它的子范畴$\mathbf{C}'$,
\begin{enumerate}
    \item $\iota$是忠实的,
    \item $\mathbf{C}'$是全子范畴当且仅当$\iota$是全的,
    \item 如果$\mathbf{C}'$是骨架, 那么$\iota$是本质满的.
\end{enumerate}

\noindent{\kaishu 证明.}

\begin{enumerate}
    \item 任取$X,Y\in\mathbf{C}'$, 显然
    $\iota\!\restriction_{\mathbf{C'}(X,Y)}:\mathbf{C'}(X,Y)\to\mathbf{C}(X,Y),
    f\mapsto f$是单射.
    \item 任取$X,Y\in\mathbf{C}'$, 显然
    $\iota\!\restriction_{\mathbf{C'}(X,Y)}:\mathbf{C'}(X,Y)\to\mathbf{C}(X,Y),
    f\mapsto f$是满射当且仅当$\mathbf{C'}(X,Y)=\mathbf{C}(X,Y)$.
    \item 如果$\mathbf{C}'$是骨架, 那么任取$X_0\in\mathbf{C}$,
    存在$X\in\mathbf{C}'$使得$\iota(X)=X\cong X_0$, 即$\iota$本质满.
    \hfill $\qedsymbol$
\end{enumerate}

\vspace{10pt}

\hypertarget{sec:例1.2}{}
\noindent\textbf{例1.2 (Hom函子)}

给定宇宙$\mathcal{U}$和$\mathcal{U}$-局部小范畴$\mathbf{C}$,
在$\mathbf{C}^{\mathrm{op}}\times\mathbf{C}$中做如下定义:
\begin{enumerate}
    \item 对任意的对象$(X,A)$, 定义$\mathrm{Hom}(X,A)=\mathbf{C}(X,A)$,
    \item 对任意的态射$(f,g):(X,A)\to(Y,B)$, 定义
    \[
        \mathrm{Hom}(f,g):\mathbf{C}(X,A)\to\mathbf{C}(Y,B),\quad
        \varphi\mapsto g\circ\varphi\circ f,
    \]
\end{enumerate}
则$\mathrm{Hom}$是从$\mathbf{C}^{\mathrm{op}}\times\mathbf{C}$到
$\mathbf{Set}$的函子, 称之为$\mathbf{C}$上的\textbf{Hom函子}.

\noindent{\kaishu 证明.}

\begin{enumerate}
    \item 因为$\mathbf{C}$是局部小范畴, 所以所有的$\mathbf{C}(X,A)$都属于$\mathcal{U}$,
    是$\mathbf{Set}$的对象.
    \item 对任意的态射$(f,g):(X,A)\to(Y,B)$和$\varphi\in\mathbf{C}(X,A)$,
    在$\mathbf{C}$中有
    \[
        f\in\mathbf{C}(Y,X),\,\,\varphi\in\mathbf{C}(X,A),\,\,g\in\mathbf{C}(A,B)
        \quad\implies\quad g\circ\varphi\circ f\in\mathbf{C}(Y,B).
    \]
    \item 对任意的对象$(X,A)$, 有
    \[
        \mathrm{Hom}(I_X,I_A):\,\varphi\,\mapsto\,
        I_A\circ\varphi\circ I_X=\varphi,\\[-6pt]
    \]
    所以$\mathrm{Hom}(I_{(X,A)})=\mathrm{Hom}(I_X,I_A)=I_{\mathrm{Hom}(X,A)}$.
    \item 对任意的态射$(f,g),(f',g')$, 在有定义的情形下,
    \vspace{3pt}
    \[\left\{\begin{aligned}
        &\mathrm{Hom}(f'\circ^{\mathrm{op}}f,g'\circ g):\quad&&
        \varphi\mapsto(g'\circ g)\circ\varphi\circ(f'\circ^{\mathrm{op}}f)
        =g'\circ g\circ\varphi\circ f\circ f',\\[-3pt]
        &\mathrm{Hom}(f',g')\circ\mathrm{Hom}(f,g):\quad&&
        \varphi\mapsto g'\circ(g\circ\varphi\circ f)\circ f'
        =g'\circ g\circ\varphi\circ f\circ f',
    \end{aligned}\right.\\[3pt]\]
    所以$\mathrm{Hom}((f',g')\circ(f,g))=\mathrm{Hom}(f'\circ^{\mathrm{op}}f,
    g'\circ g)=\mathrm{Hom}(f',g')\circ\mathrm{Hom}(f,g)$. \hfill $\qedsymbol$
\end{enumerate}

\begin{center}\begin{tikzpicture}[
    arrow/.style={-Stealth, thick}]
    \node (Cat_C) at (-5, 0.8) {$\mathbf{C}$};
    \node (Cat_Set) at (-5, -1) {$\mathbf{Set}$};
    \node (X) at (-3, 1.6) {$X$};
    \node (Y) at (0, 1.6) {$Y$};
    \node (Z) at (3, 1.6) {$Z$};
    \node (A) at (-3, 0) {$A$};
    \node (B) at (0, 0) {$B$};
    \node (C) at (3, 0) {$C$};
    \draw[arrow] (X) -- node[right] {$\varphi$} (A);
    \draw[arrow] (Y) -- node[right] {$g\varphi f$} (B);
    \draw[arrow] (Z) -- node[right] {$g'g\varphi ff'$} (C);
    \draw[arrow] (Z) --node[above] {$f'$} (Y);
    \draw[arrow] (Y) --node[above] {$f$} (X);
    \draw[arrow] (A) --node[above] {$g$} (B);
    \draw[arrow] (B) --node[above] {$g'$} (C);
    \draw[arrow] (-2.9, -1) --node[below] {$\mathrm{Hom}(f,g)$} (-0.1, -1);
    \draw[arrow] (0.1, -1) --node[below] {$\mathrm{Hom}(f',g')$} (2.9, -1);
    \draw[gray, dashed] (-5.5, -0.5) -- (4.5, -0.5);
\end{tikzpicture}\end{center}





\newpage

\section{Cat范畴}

% 这两天作息烂完, 困死了.
% 写不动了. 晚安.
% 2026.06.19

% 在清华的最后一个晚上. 再见了.
% 今当远离, 临表涕零, 不知所言.

% 自强不息  厚德载物
% 2026.06.29

\hypertarget{sec:定义2.1}{}
\noindent\textbf{定义2.1 (函子的复合)}

给定范畴$\mathbf{C},\mathbf{D},\mathbf{E}$和函子
$F:\mathbf{C}\to\mathbf{D},\,G:\mathbf{D}\to\mathbf{E}$, 令
\[
    G\circ F\overset{\mathrm{def}}{=}(G_O\circ F_O,\,G_M\circ F_M),
\]
它是从$\mathbf{C}$到$\mathbf{E}$的函子.

\vspace{10pt}

范畴和它们之间的函子可以构成$\mathbf{Cat}$范畴.

\vspace{10pt}

\hypertarget{sec:定义2.2}{}
\noindent\textbf{定义2.2 (Cat范畴)}

给定一些范畴组成的集合$O$:
\begin{enumerate}
    \item 定义$\displaystyle M=\bigcup_{\mathbf{C},\mathbf{D}\in O}
    \{(F,\mathbf{C},\mathbf{D})\mid F:\mathbf{C}\to\mathbf{D}\}$,
    \vspace{3pt}
    \item 对任意的$(F,\mathbf{C},\mathbf{D})\in M$, 定义
    $s(F,\mathbf{C},\mathbf{D})=\mathbf{C}$,
    $t(F,\mathbf{C},\mathbf{D})=\mathbf{D}$,
    \item 对任意的$\mathbf{C}\in O$, 定义$i(\mathbf{C})=I_{\mathbf{C}}$,
    \item 对任意的$(F,\mathbf{C},\mathbf{D}),(G,\mathbf{D},\mathbf{E})\in M$,
    定义$(G,\mathbf{D},\mathbf{E})\circ(F,\mathbf{C},\mathbf{D})=
    (G\circ F,\,\mathbf{C},\mathbf{E})$ (函子的复合),
\end{enumerate}
则$\mathbf{Cat}_O=(O,M,s,t,i,\circ)$是范畴.

\noindent{\kaishu 备注.
在明确$\mathbf{C},\mathbf{D}$的时候, 在不产生歧义的情况下不区分函子$F$
与态射$(F,\mathbf{C},\mathbf{D})$.}

\vspace{10pt}

\hypertarget{sec:定理2.1}{}
\noindent\textbf{定理2.1}

任给一个宇宙$\mathcal{U}$, 记$O$为所有$\mathcal{U}$-小范畴组成的集合,
则$\mathbf{Cat}_O$是$\mathcal{U}$-局部小范畴, 也简记为$\mathbf{Cat}$.

\noindent{\kaishu 证明.}

对任意的$\mathbf{C},\mathbf{D}\in O$, 有
$\mathrm{Ob}(\mathbf{C}),\mathrm{Ob}(\mathbf{D}),
\mathrm{Mor}(\mathbf{C}),\mathrm{Mor}(\mathbf{D})\in\mathcal{U}$, 所以
\[
    [\mathbf{C},\mathbf{D}]\subset
    \,\,\!^{\mathrm{Ob}(\mathbf{C})}\mathrm{Ob}(\mathbf{D})\times
    \,\,\!^{\mathrm{Mor}(\mathbf{C})}\mathrm{Mor}(\mathbf{D})\in\mathcal{U},
\]
从而
\[
    \mathbf{Cat}(\mathbf{C},\mathbf{D})=[\mathbf{C},\mathbf{D}]\times
    \{\mathbf{C}\}\times\{\mathbf{D}\}\in\mathcal{U},
\]
即$\mathbf{Cat}$是局部小范畴. \hfill $\qedsymbol$

\end{document}